% Template Oficial da Revista Physicae Organum
%   Equipe Physicae Organum
%-----------------------------

% Vá até a linha com o texto "EDITE A PARTIR DAQUI".
% Você terá que inserir título, autores, instituição, notas, resumo, palavras-chave, textos das seções e bibliografia.
% Para facilitar, colocamos o comentário "[EDITAR]" nas linhas que terão que ser alteradas.
% Indicamos alguns tutorias LaTeX ao final desta página.

% Definição do documento e pacotes utilizados. 

\documentclass[12pt,a4paper,twoside]{article}

% Configurações:

\input{conf/texto.tex}
\input{conf/matematica.tex}
\input{conf/ciencia.tex}
\input{conf/ambientes.tex}

% Para citações:

\usepackage[brazilian,hyperpageref]{backref}
\usepackage[alf]{abntex2cite}

% EDITE A PARTIR DAQUI:

\def \tituloportugues {COLOQUE SEU TÍTULO EM PORTUGUÊS AQUI, ELE PASSARÁ PARA A OUTRA LINHA AUTOMATICAMENTE SE PRECISO}  % Título. [EDITAR]
\def \tituloingles {COLOQUE O TÍTULO EM INGLÊS AQUI, ELE PASSARÁ PARA A OUTRA LINHA AUTOMATICAMENTE SE PRECISO}  % Título em inglês. [EDITAR]

% As configurações de layout estão no arquivo mostrado abaixo:
\input{conf/layout.tex} % Não mexa nesta linha!

% Coloque abaixo os nomes dos autores e o número de suas instituições:

\author[1]{Coloque Nome do Autor Um}
\author[2]{Coloque Nome do Autor Dois}
\author[1]{Coloque Nome do Autor Três}
\author[2]{Coloque Nome do Autor Quatro}

% Coloque abaixo os números e nomes das instituições:

\affil[1]{Coloque aqui o nome da instituição 1.}
\affil[2]{Coloque aqui o nome da instituição 2.}
\affil[3]{Coloque aqui o nome da instituição 3.}

% Ajustes diversos antes de começar o documento (últimos preparativos),
% tais como atribuição da variável "title", configurações para resolver erros de compilação etc.

\input{conf/ultimos-preparativos.tex}

%-----------------------------

% Aqui damos início ao documento.

\begin{document}

\maketitle

\thispagestyle{capa}

% Resumo e palavras-chave. 

\hrule

\selectlanguage{portuges}
\begin{abstract}
\noindent
Coloque aqui o texto do seu resumo. % [EDITAR] %
\vspace{\baselineskip} \\
\palavraschave{Física. História. Universidade. Modelos de universidade.}
\end{abstract}
\vspace{5mm}

\hrule

\selectlanguage{english}
\begin{abstract}
\noindent
Coloque aqui a tradução do seu resumo para o inglês. % [EDITAR] %
\vspace{\baselineskip} \\
\keywords{Physics. History. University. Models of university.}
\end{abstract}
\vspace{5mm}

\hrule

\selectlanguage{portuges}


\section{Exemplo: Introdução} % [EDITAR]
%TEXTO DA SEÇAO

Inclua o conteúdo de sua seção aqui. \\ % [EDITAR]

% Crie novas seções com o comando
% \section{Nome da seção}.

% Se preferir, crie arquivos separados com o texto de cada seção e inclua aqui assim:
% \section{Nome da seção}.
% \include{nome_do_arquivo_com_texto.tex}

\section{Exemplo: Desenvolvimento} % [EDITAR]
%TEXTO DA SEÇAO

 Inclua o conteúdo de sua seção aqui. \\ % [EDITAR]

\subsection{Exemplo: subseção} % [EDITAR]

 Inclua o conteúdo de sua subseção aqui. \\ % [EDITAR]

\section{Exemplo: Conclusão} % [EDITAR]
%TEXTO DA SEÇAO

Inclua o conteúdo de sua seção aqui. \\ % [EDITAR]

\section*{Referências} % [EDITAR]

\noindent Paardekooper S.-J. et al. (2010) Mon. Not. R. Astron. Soc., 401,
1950. % [EDITAR]

\noindent Duncan, M. J., Levison, H. F. \& Lee, M.H. (1998). A.J., 116, 2067. % [EDITAR]

\noindent Baruteau et al. (2014) in: Protoplanetary Disks and Stars VI, University of Arizona Press, 667. % [EDITAR]

\vspace{\baselineskip} \hrule

\vspace{2\baselineskip}

Se você utilizar a seção acima, apague o comando ``bibliography{main}'' do código. A referência também pode ser criada usando Bibitem, BibTeX, BibLaTeX etc. Caso você use BibTeX, apague a seção acima, insira as referências no arquivo ``main.bib''. Depois cite segundo este modelo: \cite{greenwade93}. Ou, para citar várias referências juntas, assim: \cite{greenwade93,wikimedia-latex}. Suas referências ficarão assim:

\bibliography{main}

\end{document}

% Tutoriais LaTeX:

% <http://www.mat.ufmg.br/~regi/topicos/intlat.pdf>
% <http://ece.uprm.edu/~caceros/latex/introduction.pdf>

%-----------------------------

% LaTeX Online (free account):

% <https://www.overleaf.com/>

%-----------------------------

% Códigos matemáticos LaTeX:

% <http://www.codecogs.com/latex/eqneditor.php>

%-----------------------------

% Construir tabelas no LaTeX:

% <http://www.tablesgenerator.com/#>

%-----------------------------